\section{Conclusion}

This paper introduced \benchmark{}, a benchmark for studying how agents learn from real-world environments over day-long horizons. Across 134 diverse executable tasks and roughly 38{,}000 hours of environment interaction, we find that aggregate learning trajectories follow a precise log-sigmoid relationship with interaction time. The same form appears across task families, remains stable over longer horizons, and supports forecasting later performance from early trajectories. We also find that agent learning speed has improved rapidly across recent frontier model generations.

These results suggest that learning from environments is not merely a collection of idiosyncratic task outcomes, but a measurable scaling object. Unlike many aggregate capability curves, environment-learning trajectories expose the intermediate attempts, feedback, and revisions through which progress occurs. This makes \benchmark{} useful not only for ranking agents, but for studying how agents acquire and reuse experience. More broadly, the regularity we observe suggests that post-deployment learning from rich environments may deserve the same systematic scaling attention that pretraining has received.
